\documentclass[12pt,a4paper]{article}

\usepackage[utf8]{inputenc}
\usepackage[T2A]{fontenc}
\usepackage[russian]{babel}
\usepackage{amsmath, amsfonts, amssymb, amsthm}
\usepackage{geometry}
\geometry{top=2cm, bottom=2cm, left=2.5cm, right=1.5cm}
\usepackage{tikz}
\usetikzlibrary{arrows.meta, positioning}
\usepackage{algorithm}
\usepackage{algpseudocode}
\usepackage{hyperref}

\newtheorem{theorem}{Теорема}
\newtheorem{definition}{Определение}
\newtheorem{example}{Пример}

\title{Билет №81 \\ \small Методы представления знаний: процедурные представления, логические представления, семантические сети, фреймы, системы продукций. Интегрированные методы представления знаний. Языки представления знаний. Базы знаний. (ИТМО) \\ Методы представления знаний: процедурные представления, логические представления, семантические сети, фреймы, системы продукций. Интегрированные методы представления знаний. Языки представления знаний. Базы знаний. (СпБГУ)}
\author{Сериков Владислав Александрович}
\date{16 мая 2026}

\begin{document}

\maketitle
\tableofcontents
\newpage

\section{Введение}

\begin{definition}
\textbf{База знаний (БЗ)} — это семантическая модель, описывающая предметную область и позволяющая отвечать на вопросы из этой предметной области, выводить новые факты и решать практические задачи. БЗ состоит из двух частей: экстенсионала (конкретные факты) и интенсионала (правила, закономерности, абстрактные сущности).
\end{definition}

Представление знаний является ключевой задачей при построении интеллектуальных систем. Существуют различные подходы, каждый из которых имеет свои преимущества и области применения.

\section{Логические представления}

Логическое представление основано на формальных логических системах (логика высказываний, логика предикатов первого порядка и неклассические логики). Знания описываются в виде набора аксиом (фактов и правил), а получение новых знаний сводится к формальному логическому выводу.

\begin{definition}
\textbf{Логика предикатов первого порядка (FOL)} — формальная система, включающая переменные, константы, функции, предикаты, логические связки ($\land, \lor, \neg, \to$) и кванторы ($\forall, \exists$).
\end{definition}

Основным механизмом логического вывода в автоматических системах является метод резолюций, основанный на правиле резолюции. Для логики высказываний правило резолюции выглядит так: из дизъюнктов $(A \lor p)$ и $(B \lor \neg p)$ следует $(A \lor B)$.

\begin{theorem}[О полноте метода резолюций для логики высказываний]
Если конечное множество дизъюнктов $S$ невыполнимо (то есть не имеет интерпретации, при которой все дизъюнкты истинны), то из $S$ с помощью конечного числа применений правила резолюции можно вывести пустой дизъюнкт $\square$ (символ тождественной лжи).
\end{theorem}

\begin{proof}
Доказательство проведем индукцией по числу различных пропозициональных переменных (атомов) $n$ в множестве $S$.

\textbf{База индукции.} При $n=0$ невыполнимое множество может состоять только из пустого дизъюнкта $\square$. Вывод тривиален, так как $\square$ уже содержится в $S$.

\textbf{Индуктивный переход.} Предположим, что теорема верна для $n$ атомов. Рассмотрим множество $S$ с $n+1$ атомами. Выберем произвольный атом $p$.
Определим множества:
1. $S_{p=0}$ — множество, полученное из $S$ удалением всех дизъюнктов, содержащих $\neg p$, и удалением литерала $p$ из оставшихся.
2. $S_{p=1}$ — множество, полученное из $S$ удалением всех дизъюнктов, содержащих $p$, и удалением литерала $\neg p$ из оставшихся.

Поскольку $S$ невыполнимо, то и $S_{p=0}$, и $S_{p=1}$ также невыполнимы (иначе выполняющая оценка для них дала бы оценку для $S$). В обоих множествах не более $n$ атомов.
По предположению индукции, из $S_{p=0}$ и из $S_{p=1}$ можно вывести $\square$.
Если мы повторим шаги вывода $\square$ из $S_{p=0}$ на исходном множестве $S$, мы получим либо $\square$ (если литерал $p$ не участвовал), либо дизъюнкт $\{p\}$. Аналогично, из $S_{p=1}$ вывод на исходном множестве даст либо $\square$, либо $\{\neg p\}$.

Если хотя бы один вывод дает $\square$, теорема доказана. Если же получены дизъюнкты $\{p\}$ и $\{\neg p\}$, применим к ним правило резолюции и получим пустой дизъюнкт $\square$. Переход доказан.
\end{proof}

\begin{algorithm}
\caption{Алгоритм вывода методом резолюций (опровержение)}\label{alg:resolution}
\begin{algorithmic}[1]
\State \textbf{Вход:} База знаний $KB$, целевое утверждение $\alpha$.
\State \textbf{Выход:} \texttt{True}, если $KB \models \alpha$, иначе \texttt{False}.
\State Сформировать множество формул $S = KB \cup \{\neg \alpha\}$.
\State Привести все формулы в $S$ к конъюнктивной нормальной форме (КНФ).
\Loop
    \State Найти в $S$ два дизъюнкта $C_1$ и $C_2$, содержащие контрарные литералы ($L$ и $\neg L$).
    \State Сформировать резольвенту $R$ (дизъюнкция $C_1$ и $C_2$ без $L$ и $\neg L$).
    \If{$R == \square$ (пустой дизъюнкт)}
        \State \Return \texttt{True}
    \EndIf
    \If{$R$ уже содержится в $S$}
        \State \textbf{continue}
    \EndIf
    \State Добавить $R$ в $S$.
    \If{новых резольвент не возникает}
        \State \Return \texttt{False}
    \EndIf
\EndLoop
\end{algorithmic}
\end{algorithm}

\begin{example}
\textbf{KB:} Человек смертен ($H \to M$). Сократ — человек ($H$).\\
\textbf{Цель:} Сократ смертен ($M$).\\
\textbf{Опровержение:} $S = \{\neg H \lor M, H, \neg M\}$. \\
Резолюция 1: Из $H$ и $\neg H \lor M$ получаем $M$.\\
Резолюция 2: Из $M$ и $\neg M$ получаем $\square$. \textbf{Вывод успешен.}
\end{example}

\section{Процедурные представления}

В процедурных представлениях знания кодируются непосредственно в виде исполняемого кода (алгоритмов, подпрограмм). В отличие от декларативных моделей, где знания отделены от механизма вывода, здесь знания и способы их использования слиты воедино.

\textbf{Преимущества:} высокая скорость исполнения, возможность моделирования сложных динамических процессов.\\
\textbf{Недостатки:} сложность обновления и модификации базы знаний, "непрозрачность" для пользователя (сложно объяснить причину принятия того или иного решения).

Примером процедурного представления служат демоны во фреймовых системах или триггеры в реляционных базах данных.

\section{Системы продукций}

Система продукций (правил) состоит из трех основных компонентов:
1. \textbf{Рабочая память (РП)} — хранит текущие факты.
2. \textbf{База правил (БП)} — содержит набор правил вида \texttt{ЕСЛИ <условие> ТО <действие>}.
3. \textbf{Машина вывода (МВ)} — механизм применения правил к фактам.

Выделяют прямой вывод (от фактов к цели) и обратный вывод (от цели к фактам).

\begin{algorithm}
\caption{Алгоритм прямого вывода}\label{alg:forward_chaining}
\begin{algorithmic}[1]
\State \textbf{Инициализация:} Поместить исходные факты в Рабочую Память (РП).
\Loop
    \State \textbf{Сопоставление:} Найти конфликтное множество — все правила из БП, условия которых истинны при текущем состоянии РП.
    \If{конфликтное множество пусто}
        \State \textbf{Остановка:} Выход из цикла.
    \EndIf
    \State \textbf{Разрешение конфликта:} Выбрать одно правило из конфликтного множества согласно стратегии (например, правило с наиболее специфичным условием или первое по порядку).
    \State \textbf{Действие:} Выполнить правую часть выбранного правила (обычно это добавление новых фактов в РП).
\EndLoop
\end{algorithmic}
\end{algorithm}

\section{Семантические сети}

Семантическая сеть — это ориентированный граф, вершины которого соответствуют объектам, понятиям или ситуациям, а дуги — отношениям между ними.

Основные типы отношений:
\begin{itemize}
    \item \textbf{is-a} (является экземпляром/подклассом) — формирует иерархию и обеспечивает наследование свойств.
    \item \textbf{part-of} (является частью).
    \item Атрибутивные отношения (имеет цвет, размер и т.д.).
\end{itemize}

\begin{center}
\begin{tikzpicture}[node distance=2cm, every node/.style={draw, rectangle, rounded corners, fill=blue!10, minimum width=2.5cm}, >=Latex]
    \node (animal) {Животное};
    \node (bird) [below=of animal] {Птица};
    \node (penguin) [below left=of bird] {Пингвин};
    \node (canary) [below right=of bird] {Канарейка};
    \node (wings) [right=of bird] {Крылья};

    \draw[->] (bird) -- node[auto, draw=none, fill=none] {is-a} (animal);
    \draw[->] (penguin) -- node[auto, draw=none, fill=none] {is-a} (bird);
    \draw[->] (canary) -- node[auto, draw=none, fill=none] {is-a} (bird);
    \draw[->] (bird) -- node[auto, draw=none, fill=none] {has-part} (wings);
    \draw[->] (penguin) -- node[auto, draw=none, fill=none, sloped] {can-fly (False)} (wings); % Exception
\end{tikzpicture}
\end{center}

В семантических сетях вывод сводится к операциям поиска на графе с учетом наследования (например, Канарейка имеет крылья, так как она Птица).

\section{Фреймы}

Концепция фреймов предложена М. Минским. \textbf{Фрейм} — это структура данных для представления стереотипной ситуации или объекта. Фрейм состоит из слотов (атрибутов).

Каждый слот может содержать:
\begin{itemize}
    \item Конкретное значение (например, \texttt{Цвет: Красный}).
    \item Значение по умолчанию (default).
    \item Процедуры (демоны): \texttt{IF-NEEDED} (запускается при запросе значения), \texttt{IF-ADDED} (при изменении значения).
    \item Указатель на другой фрейм (реализация отношений).
\end{itemize}

\begin{example}
\textbf{Фрейм:} Студент \\
\textbf{IS-A:} Человек \\
\textbf{Слоты:}
\begin{itemize}
    \item ВУЗ: (по умолчанию) Неизвестен
    \item Курс: (тип) Целое число от 1 до 6
    \item Средний балл: (IF-NEEDED) вычислить\_балл()
\end{itemize}
\end{example}

\section{Интегрированные методы представления знаний}

На практике ни один из базовых методов не является универсальным. \textbf{Интегрированные (гибридные) методы} объединяют несколько парадигм для компенсации недостатков одних методов достоинствами других.

Примеры интеграции:
\begin{itemize}
    \item \textbf{Продукционно-фреймовые системы:} Объекты предметной области и их структура описываются фреймами, а динамическое поведение и логика принятия решений — правилами (продукциями). Например, в правиле условие может проверять значение слота фрейма.
    \item \textbf{Логико-семантические системы:} Использование дескриптивных логик (строгой математической базы) поверх семантических сетей для гарантии корректности и разрешимости вывода.
\end{itemize}

\section{Языки представления знаний и базы знаний}

Для реализации описанных методов созданы специализированные языки:
\begin{itemize}
    \item \textbf{Prolog} — язык логического программирования, основанный на усеченном логическом выводе (Хорновские дизъюнкты) и механизме обратного вывода (унификация и поиск с возвратом).
    \item \textbf{LISP} — исторически первый язык для ИИ, отлично подходящий для процедурного и списочного представления.
    \item \textbf{CLIPS / JESS} — языки для создания продукционных экспертных систем, реализующие быстрый алгоритм сопоставления с образцом (алгоритм Rete).
    \item \textbf{OWL, RDF} — современные языки на базе XML для семантической паутины (Semantic Web), основанные на дескриптивных логиках для описания онтологий.
\end{itemize}

Базы знаний, построенные на этих языках, являются ядром экспертных систем (ЭС). Классическая архитектура ЭС включает в себя: базу знаний, машину вывода, подсистему объяснений, рабочую память и интерфейс инженера по знаниям/пользователя. Главное отличие БЗ от классических баз данных (БД) — способность генерировать новые знания из существующих с использованием механизмов вывода.

\end{document}
